\chapter{Discussão}
\label{cap:discussao}

\notaguia{Primeira redação, a partir do quadro de posicionamento do
Capítulo~\ref{cap:referencial}, dos resultados do
Capítulo~\ref{cap:resultados} e de
\texttt{Textos/referencia/referencial\_marcha.md}. Vale para este capítulo a
mesma restrição bibliográfica declarada no Capítulo~\ref{cap:referencial}:
autores mencionados sem vínculo em \texttt{ref/referencias.bib} aguardam
conferência de autoria e edição.}

Os resultados do capítulo anterior descrevem um estado que se reorganizou
sem mudar de tamanho: a mesma superfície, ocupada de outro modo, com a
conversão migrando do Sul para o Norte ao longo de quatro décadas. Este
capítulo pergunta o que essa descrição acrescenta ao que já se sabia, o que
ela implica para o Cerrado que resta e o que, do que foi apresentado,
sustenta peso de conclusão e o que apenas sustenta peso de indício.

A ordem segue a do referencial: o posicionamento na literatura --- a fronteira
e a renda da terra, o deslocamento indireto, o motor macroeconômico e o
espelho do Cerrado ---, as implicações ambientais e de política, e a
consolidação do alcance e dos limites, que até aqui apareceram dispersos junto
de cada resultado.

\section{A fronteira, a renda da terra e uma transição truncada}
\label{sec:disc-fronteira}

O achado que melhor dialoga com o referencial não é a marcha em si, e sim a
sua forma: camadas que sobem juntas sem que o vão entre elas se feche em
quatro décadas, e com magnitudes de deslocamento próximas. É a assinatura que o modelo de localização de von Thünen prevê
quando a ordenação por renda da terra é estável e o que se desloca é o
conjunto: não uma linha que avança, mas um arranjo que se translada sobre um
gradiente.

A contribuição de método por trás desse resultado vale enunciar, porque a
alternativa era circular. A ordenação ricardiana costuma ser \emph{assumida}
e depois ilustrada pelo próprio uso da terra que se quer explicar; aqui o
gradiente, medido por um dado físico exógeno, reproduziu sozinho a ordenação
Sul--Centro--Norte (Seção~\ref{sec:res-perna3}). A premissa do arcabouço
tornou-se uma medida independente dele.

Quanto à síntese de \citeonline{Angelsen2007}, a previsão era que gradiente
espacial e pico temporal fossem faces do mesmo processo. É o que se observa:
as duas aparecem simultaneamente e em lugares diferentes do mesmo estado ---
o gradiente no eixo Sul--Norte, o pico na região que abriu primeiro. Poder
observá-las sob um único desenho analítico, na mesma malha e com o mesmo
classificador, é o que o recorte estadual oferece e o recorte de bioma não.

Já a transição florestal tem uma das duas metades de sua previsão contrariada.
O Sul goiano cumpre a primeira --- estoque reduzido, taxa de conversão
cadente, estabilização a partir de 2019 --- e não cumpre a segunda: não há
regeneração líquida. A estabilização é depleção, não recuperação: perde-se
pouco porque restou pouco a perder, não porque a vegetação esteja voltando.

Essa leitura é o que sobra depois de uma verificação específica. Nas matrizes
de transição existe um fluxo reverso
volumoso de pastagem para vegetação natural, que seria a evidência natural de
regeneração. Examinado em classe bruta, ele se mostrou dirigido quase
inteiramente à formação savânica e aproximadamente simétrico em relação ao
fluxo de ida --- assinatura de oscilação do classificador na borda entre
pasto e cerrado, não de recuperação
(Seção~\ref{sec:met-decisoes}). O componente real de regeneração é
minoritário, restrito à formação florestal, e só se acumula em horizontes
longos. A implicação para o referencial é direta: o caso goiano sugere que a
fase tardia da transição florestal pode consistir apenas em exaustão da
oferta, sem a virada de sinal que o modelo associa a ela --- uma transição
truncada. Se isso é peculiaridade do Cerrado, do período ou do estado, este
desenho não decide; mas registra que a estabilização observada não autoriza a
leitura otimista que o arcabouço sugere.

Quanto à tradição brasileira, a distinção de \citeonline{Martins1996} entre
frente de expansão e frente pioneira dá o vocabulário mais preciso disponível
para os ``dois Goiáses'', com uma calibragem: as categorias de Martins são
sociológicas --- relações de trabalho, formas de propriedade, grau de mediação
capitalista da ocupação --- e o que este trabalho mede são coberturas do solo
e agregados econômicos municipais. A correspondência é analógica e serve para
nomear o contraste; não constitui teste da tese de Martins, e nada nos dados
aqui reunidos poderia constituí-lo. O mesmo vale para a leitura geopolítica de
Becker: informa a interpretação do Norte goiano como território em disputa,
mas o desenho não observa Estado, projeto ou agente.

\section{O deslocamento indireto: um nulo onde havia razão para esperar efeito}
\label{sec:disc-iluc}

O resultado de maior consequência para o posicionamento é um nulo, e a forma
de enunciá-lo determina se ele contribui ou se apenas confunde.

O canal definido na Seção~\ref{sec:ref-iluc} --- empurrão de curta distância
e curta defasagem, de uma metade do estado sobre a outra --- não apresentou,
em trinta e seis recortes, a assinatura que exigiria
(Seção~\ref{sec:res-perna3}). Como a Seção~\ref{sec:ref-iluc} registrou, esse
nulo não é previsão cumprida: havia razão teórica para esperar o efeito, e o
precedente mais próximo, no mesmo nível municipal, encontrou-o na Amazônia.
Como diferem a região, o desfecho medido e a forma da ligação espacial, a
divergência em relação a \citeonline{Arima2011} delimita o que resta comparar
em vez de estabelecer contradição.

Duas precisões impedem que essa leitura cresça além do que sustenta. A
primeira já foi enunciada e se repete aqui porque é a mais fácil de perder na
citação de terceiros: o trabalho não refuta o deslocamento indireto de uso da
terra. Ausência da assinatura testada não é ausência do fenômeno, e ficam
expressamente fora do alcance do teste um deslocamento de longo alcance que
pule os vizinhos, um de defasagem longa, um efeito temporal moderado que o
poder do teste não alcança e qualquer canal que atravesse a divisa do estado.
Ficam descartados, com margem confortável, um empurrão local visível na série
e um empurrão temporal forte.
A segunda é que o veredito temporal é \emph{sem líder}, e não ``o Sul não
lidera, logo o Norte lidera''. O método adequado a séries integradas zera a
precedência nas duas direções, e o episódio que produziu essa correção ---
uma precedência forte no sentido inverso, que se revelou artefato de
integração --- está registrado como resultado superado justamente para que a
simetria do veredito fique visível.

O que ocupa o lugar do empurrão é substancial e positivo. O maior coeficiente
de toda a frente é o da substituição local: dentro do próprio município,
quando a lavoura cresce, o pasto recua. É o desfecho observado da disputa por
terra que \citeonline{Cohn2014} projetam em modelo: no Sul goiano ela aparece
resolvida, e resolvida a favor da lavoura.
A literatura de intensificação e a de
deslocamento costumam ser mobilizadas em conjunto, como se o mesmo dado
sustentasse as duas; aqui elas se separam empiricamente: o mecanismo local
aparece em quase toda especificação, e o mecanismo à distância não aparece em
nenhuma.

Um detalhe do resultado espacial pede prudência interpretativa. As doze
estimativas são todas negativas --- isto é, quando a lavoura dos vizinhos ao
sul cresce, o pasto local diminui ---, e onze delas não são significativas;
a décima segunda é significativa, mas no sentido \emph{oposto} ao que a
hipótese de empurrão exigia. A leitura
tentadora seria transformar isso num achado sobre transbordamento invertido.
Ela não se sustenta com este desenho: o mais que se pode dizer é que a
substituição local tem alguma expressão espacial --- municípios vizinhos
compartilham a mesma dinâmica de troca de pasto por lavoura --- e que isso é
o oposto do que a hipótese do empurrão exigia. O sinal é informativo por
excluir a hipótese testada, não por estabelecer outra em seu lugar.

\section{O motor comum e a divisão de trabalho com a literatura}
\label{sec:disc-drive}

Se as duas metades do estado não se comandam, resta explicar por que marcham
juntas, e a resposta proposta --- resposta coordenada a um estímulo externo
comum, incidindo sobre um gradiente de aptidão --- precisa ser apresentada
com o peso que a evidência lhe dá, que não é o peso de uma demonstração.

A divisão de trabalho com a literatura, enunciada na
Seção~\ref{sec:ref-cambio}, é o que torna a leitura defensável: o nível
causal do canal cambial vem de \citeonline{Richards2012}, em escala nacional
e continental, e o que esta pesquisa estima é apenas a geografia desse
efeito dentro do estado.

Sobre esse gradiente, o veredito é corroborante e não estabelecido. O
coeficiente tem a direção prevista, mas o p-valor apropriado ao desenho ---
permutação do \emph{shifter}, que substitui um erro-padrão agrupado
reconhecidamente otimista para desenhos com um único choque --- fica entre
0,07 e 0,13, sem cruzar o corte convencional. Os p-valores agrupados de 0,026
e 0,031, produzidos pelo desenho anterior, não devem ser citados como
significância. O que sustenta a leitura não é
a significância, e sim a especificidade: os placebos de desfecho saem vazios,
os placebos no tempo saem vazios e o sinal é estável no \emph{jackknife} por
ano. Um padrão que aparece só onde deveria aparecer é evidência de tipo
diferente da que um p-valor fornece, e menos conclusiva do que ela.

O teto dessa frente é estrutural, e convém dizer o que ele exigiria para ser
levantado. Um choque
macroeconômico nacional oferece da ordem de 38 realizações independentes, e
nenhuma quantidade adicional de unidades espaciais aumenta esse número: o
poder é limitado pela dimensão temporal. Levantar o teto exigiria um
\emph{shifter} com variação transversal própria --- exposição diferencial a
preços por produto, por exemplo, em vez de uma única série nacional --- ou
uma janela temporal maior, ou um painel de vários estados que multiplicasse
as realizações do choque. São desenhos distintos, não robustez adicional
sobre este.

Segue disso uma delimitação que precisa acompanhar qualquer citação deste
resultado: o que a pesquisa estabelece é o \emph{mecanismo} --- reorganização
sob força comum, e não empurrão entre regiões ---, e não a identidade da
força. O câmbio é o candidato mais forte e permanece candidato. Qual variável
exatamente aciona o gatilho, o dado estadual não tem como cravar, porque
preço internacional, câmbio, crédito e expectativa de renda movem-se juntos e
são absorvidos pelo mesmo efeito fixo de ano.

\section{Goiás como microcosmo do Cerrado}
\label{sec:disc-matopiba}

A previsão enunciada na Seção~\ref{sec:ref-matopiba} era a de que o contraste
que \citeonline{CarneiroFilho2016} descrevem no bioma --- expansão sobre
pastagem já aberta no Cerrado consolidado, sobre vegetação natural no
MATOPIBA --- reaparecesse dentro de Goiás. Ela se cumpre: o Sul comporta-se
como o Cerrado maduro, com expansão sobre pasto e a substituição local como
mecanismo dominante, e o Norte como o MATOPIBA, com a abertura de vegetação
natural sustentando a fronteira. Os números que separam os dois regimes são
os que menos dependem da classe ambígua, porque nem a origem nem o destino
passam por ela, e a assimetria entre as pontas do estado é de quase quatro
para um na queda dessa conversão entre os Atos~II e III.

O ganho dessa coincidência é metodológico antes de ser substantivo. Comparar
regimes de fronteira entre regiões do bioma implica comparar, junto com eles,
convenções de medida diferentes: recortes distintos, séries estatísticas de
cobertura desigual, malhas administrativas com histórias próprias. Dentro de
um mesmo estado, os dois regimes são observados sob o mesmo classificador, a
mesma malha de território constante, a mesma fonte estatística e o mesmo
conjunto de decisões metodológicas. O que separa Sul e Norte não pode ser
atribuído à régua, porque a régua é uma só.

Uma tentação deve ser resistida: ler o Sul goiano como o futuro do
MATOPIBA. Os dois estão em
posições diferentes de um processo semelhante, mas nada nos dados garante
que a trajetória se repita --- estrutura fundiária, infraestrutura, momento
tecnológico e cobertura de proteção diferem, e a própria transição truncada
descrita na Seção~\ref{sec:disc-fronteira} mostra que o fim da trajetória não
é dado pelo arcabouço. O que a comparação autoriza é mais modesto e ainda
assim útil: o Sul mostra o que acontece com a fronteira quando a fonte de
terra passa de vegetação natural para pasto já aberto, e mostra que essa
passagem pode ocorrer com a demanda no pico.

\section{Implicações ambientais}
\label{sec:disc-ambiental}

A primeira implicação é de exposição, e o vocabulário importa. Restam
6,56~Mha de vegetação natural ainda exposta à conversão --- savana e campo,
as classes que a fronteira de fato consome, e não a totalidade dos
11,88~Mha de vegetação natural em pé em 2024 ---, cerca de 60\% do que havia
em 1985 na mesma definição, e o Norte
concentra 44\% desse total --- precisamente a região para onde a conversão
migrou e onde a taxa não caiu. Esse número mede o Cerrado que ainda pode ser
perdido; ele não descreve uma reserva de terra aproveitável, e tratá-lo assim
inverteria o sentido do resultado. É também um teto, por razão declarada: sem
cadastro ambiental integrado pixel a pixel, o estoque inclui Reserva Legal e
APP que na prática não podem ser convertidas, de
modo que a área realmente exposta é menor que 6,56~Mha, nunca maior.

A segunda implicação é sobre onde a proteção está. Os números da
Seção~\ref{sec:res-perna4} --- proteção integral abaixo de 3\% em qualquer
região, congelada desde 1985, crescendo 0,12~Mha contra 4,11~Mha suprimidos
--- não sustentam a leitura de que a lei barrou a fronteira do Sul. A
conclusão a extrair não é que a proteção falhou onde tentou: é que ela não
estava no caminho, e não está agora, já que o Norte concentra 44\% do que
resta exposto sem ter pela frente área de proteção integral que o barre.
Reserva Legal e APP são outra questão, e permanecem em aberto pelas razões
ali declaradas.

A terceira implicação vem da contabilidade de carbono e é a menos intuitiva.
Como o custo de carbono não é proporcional à área --- a formação florestal
perde bem menos área que a savânica e ainda assim emite mais ---, uma métrica
de política formulada em hectares, seja de área convertida, protegida ou
recuperada, precifica mal o carbono: trata como equivalentes hectares que não
o são. Some-se a isso a distribuição temporal: com 80\% da emissão saindo
antes de 2001, o grosso do dano climático da marcha antecede o período em que
a marcha passou a ser observada, o que desloca a pergunta de mitigação do
passivo acumulado --- irreversível na escala de tempo relevante --- para a
exposição que resta e para onde ela está.

A quarta implicação é social e exige a maior contenção das quatro. A
fronteira norte quase dobrou a área cultivada sem que o vão de
desenvolvimento municipal em relação ao Sul se fechasse. A leitura é
associativa, com o gradiente espacial controlado, e não causal; o IFDM é
índice composto e municipal, e não capta distribuição interna. Com essas
ressalvas, o que se pode dizer é que a expansão da área cultivada, no período
e na escala observados, não aparece associada a convergência de
desenvolvimento --- o que recomenda cautela com o argumento de que abrir área
é, por si, política de desenvolvimento regional. Estabelecer o contrário
exigiria desenho que este trabalho não tem.

\section{Implicações para a política pública e para o monitoramento}
\label{sec:disc-politica}

A distinção entre empurrão e motor comum não é escolástica: ela muda a
alavanca. Se a conversão no Norte fosse efeito-dominó da lavoura do Sul,
conter a lavoura do Sul conteria a fronteira, e a política eficaz seria de
realocação. Na medida em que as duas dinâmicas respondam a um estímulo comum
sobre um gradiente de aptidão --- leitura corroborante, não estabelecida
---, uma política que apenas desloque
atividade entre regiões tende a não reduzir a conversão total: o que muda o
resultado agregado é o que altera o retorno de converter, e isso age sobre o
estado inteiro.

Essa implicação precisa ser lida junto com um limite de desenho declarado no
Capítulo~\ref{cap:metodologia} e que o trabalho não contorna. Os marcos
institucionais examinados são federais e não deixam grupo não tratado, de
modo que nenhuma estimativa aqui identifica efeito de política. O trabalho não
diz, portanto, qual instrumento funciona. O que ele diz é geográfico e
temporal, e isso já é acionável: onde o instrumento teria de estar --- à
frente da fronteira, no Norte --- e quando --- antes de o fluxo chegar,
porque a proteção que existe veio depois da conversão e está ao sul dela.

O teste de simetria reforça o ponto por outro caminho. Nenhuma das camadas
que poderiam puxar a fronteira aparece à frente dela: o crédito rural fica
cerca de 70~km ao sul da pastagem, e a capacidade de armazenagem, cerca de
150~km, praticamente sobre o núcleo de lavoura do Sudoeste. Silo, banco e
cadeia exportadora chegam depois, não antes. Nada disso mede capacidade de
puxar --- é posição, não efeito ---, mas contraria a suposição, comum no debate, de que restringir
infraestrutura ou crédito seja o ponto de aplicação natural para deter uma
fronteira que já se moveu.

Há, por fim, uma implicação operacional para quem monitora. No trecho final da série, a
conversão de pastagem em agricultura passa a ser crescentemente rotulada como
``Mosaico de Usos''. Qualquer indicador, meta ou gatilho de política que se
apoie na classe ``Agricultura'' da série anual subconta a lavoura recente
exatamente na janela em que a decisão precisa ser tomada --- e subconta de
forma desigual no território, mais ao norte do que ao sul, porque é lá que a
conversão nova está. A leitura errada disponível é confortável: a de que a
conversão desacelerou. A resposta adequada não é somar as duas classes, o que
superconta na mesma medida em que a classe isolada subconta, e sim reportar o
intervalo entre as duas réguas e ancorá-lo numa fonte que não passe pelo
classificador --- área plantada, rebanho, crédito. É recomendação de
procedimento, de custo baixo e transferível a qualquer análise que use a
mesma série.

\section{Alcance e limites do que foi estabelecido}
\label{sec:disc-limites}

Os limites de cada resultado foram declarados junto dele, no
Capítulo~\ref{cap:resultados}. Consolidá-los torna visível o alcance conjunto
do que se afirma e separa o que já não melhora com mais teste do que melhora
com outro desenho (Quadro~\ref{quadro:alcance}).

\begin{quadro}[htb]
\centering
\caption[Grau de estabelecimento por frente]{Grau de estabelecimento por frente e o que elevaria cada um.}
\label{quadro:alcance}
\small
\begin{tabular}{p{2.6cm}p{3.6cm}p{2.4cm}p{5.0cm}}
\toprule
\textbf{Frente} & \textbf{Afirmação} & \textbf{Grau} & \textbf{O que elevaria} \\
\midrule
1. O padrão &
A conversão e o rebanho deslocam-se ao norte; o gradiente latitudinal
persiste &
estabelecido &
nada a elevar: sobrevive a pixel, malha, classe e barra de erro \\
\addlinespace
2. Mecanismo &
Duas lógicas coexistem em todo o estado; a região explica pouco da idade da
pastagem &
estabelecido &
o que a região não explica pede dado por imóvel rural, não agregado
municipal \\
\addlinespace
3a. Deslocamento &
O canal local e contemporâneo não apresenta a assinatura exigida &
nulo com poder declarado &
desenho de maior alcance espacial e defasagem longa; travessia da divisa
estadual \\
\addlinespace
3b. \emph{Drive} comum &
O choque comum desloca o rebanho para onde a aptidão é menor &
corroborante &
\emph{shifter} com variação transversal, janela maior ou painel
multiestadual \\
\addlinespace
4. Teto de oferta &
A freada do Sul é de oferta de terra, não de demanda nem de proteção &
sustentado por eliminação &
cadastro ambiental pixel a pixel; e o Ato~III é curto \\
\bottomrule
\end{tabular}
\fontefig{elaboração própria.}
\end{quadro}

Alguns limites atravessam todas as frentes e merecem enunciado próprio.

\textbf{O alcance é intra-estadual.} Tudo o que foi testado ocorre dentro de
Goiás. Soja goiana deslocando pecuária para o Pará ou para o Maranhão é outra
pergunta, e os dados aqui reunidos não a fazem. Nenhuma afirmação deste
trabalho sobre deslocamento indireto se estende para fora da divisa.

\textbf{A medida depende de um classificador, e ele mudou de comportamento no
fim da série.} O intervalo entre as duas réguas contém o artefato; ele não o
corrige. Converter o intervalo em ponto exigiria discriminar a reetiquetagem
dos sistemas integrados reais, o que demandaria comparação com a coleção
anterior do MapBiomas --- trabalho desenhado e não executado. Enquanto isso
não for feito, toda medida de agricultura em janela que alcance 2020--2024
permanece um intervalo, e as conclusões que dependem da convenção de classe
permanecem declaradas como tais.

\textbf{O terceiro ato tem cinco anos.} A freada do Sul, a mudança de endereço
da fronteira e a estabilização do estoque meridional são sinal inicial, não
regime consolidado. A estabilização pode ainda revelar-se uma pausa, e o
desenho não distingue as duas hipóteses --- só o tempo o fará.

\textbf{``Convertível'' e ``protegida'' são aproximações declaradas.} A
primeira é o Cerrado que o MapBiomas ainda vê em pé; a segunda, a malha
oficial de unidades de conservação. São tetos adequados ao argumento de
estoque, mas medem exposição por aproximação, não lote a lote.

\textbf{O recorte por mesorregião é grosso.} Cinco faixas para um estado do
tamanho de Goiás escondem variação interna; onde foi possível descer à malha
fina --- as 166 AMCs, e até o pixel ---, a conclusão não mudou.

\textbf{Sobre a idade da pastagem, nada se afirma quanto a tendência.} O eixo
do tempo está comprometido nas duas pontas: antes, pela idade truncada no
início da série; depois, pela mudança de rótulo. Que o primeiro período
apareça com uma população só não significa que a reserva antiga tenha surgido
depois, e que os dois grupos apareçam mais separados no período recente não
significa que estejam se separando. Nos dois casos, o que se lê é o tamanho da
janela de observação.

A esses soma-se um limite de resolução que a própria pesquisa encontrou e não
resolveu, e que talvez seja o achado negativo mais útil para quem continuar.
A escolha entre girar o capim em poucos anos e deixá-lo por décadas resistiu
a todos os preditores agregados testados, do recorte regional ao perfil
produtivo e ao crédito. Ela acontece abaixo da escala que dados agregados
municipais alcançam --- no talhão, na propriedade, na história daquela fazenda
--- e a consequência é que a próxima pergunta desta linha de investigação
exige outra unidade de observação, não mais tratamento estatístico da mesma.
